
\documentclass{article}
\usepackage{colortbl}
\usepackage{makecell}
\usepackage{multirow}
\usepackage{supertabular}

\begin{document}

\newcounter{utterance}

\centering \large Interaction Transcript for game `bbh', experiment `logical\_deduction\_seven\_objects', episode 55 with SLED{-}BSU\_\_Qwen3.5{-}27B{-}sft{-}ep1\_\_prm{-}ep1{-}best\_of\_n.
\vspace{24pt}

{ \footnotesize  \setcounter{utterance}{1}
\setlength{\tabcolsep}{0pt}
\begin{supertabular}{c@{$\;$}|p{.15\linewidth}@{}p{.15\linewidth}p{.15\linewidth}p{.15\linewidth}p{.15\linewidth}p{.15\linewidth}}
    \# & \multicolumn{2}{c}{Player} && \multicolumn{2}{c}{Game Master} \\
    \hline

    \theutterance \stepcounter{utterance}  
    & & & \multicolumn{4}{p{0.6\linewidth}}{
        \cellcolor[rgb]{0.9,0.9,0.9}{
            \makecell[{{p{\linewidth}}}]{
                \texttt{\tiny{[P1$\langle$GM]}}
                \texttt{A logical deduction task which requires deducing the order of a sequence of objects.} \\
\\ 
\texttt{Q: The following paragraphs each describe a set of three objects arranged in a fixed order. The statements are logically consistent within each paragraph. In a golf tournament, there were three golfers: Amy, Eli, and Eve. Eve finished above Amy. Eli finished below Amy.} \\
\texttt{Options:} \\
\texttt{(A) Amy finished last} \\
\texttt{(B) Eli finished last} \\
\texttt{(C) Eve finished last} \\
\texttt{A: (B)} \\
\\ 
\texttt{Q: The following paragraphs each describe a set of three objects arranged in a fixed order. The statements are logically consistent within each paragraph. On a shelf, there are three books: a white book, a green book, and an orange book. The green book is to the right of the white book. The orange book is the rightmost.} \\
\texttt{Options:} \\
\texttt{(A) The white book is the leftmost} \\
\texttt{(B) The green book is the leftmost} \\
\texttt{(C) The orange book is the leftmost} \\
\texttt{A: (A)} \\
\\ 
\texttt{Q: The following paragraphs each describe a set of three objects arranged in a fixed order. The statements are logically consistent within each paragraph. On a shelf, there are three books: a red book, a gray book, and a white book. The white book is to the left of the gray book. The red book is the second from the left.} \\
\texttt{Options:} \\
\texttt{(A) The red book is the leftmost} \\
\texttt{(B) The gray book is the leftmost} \\
\texttt{(C) The white book is the leftmost} \\
\texttt{A: (C)} \\
\\ 
\texttt{Q: The following paragraphs each describe a set of seven objects arranged in a fixed order. The statements are logically consistent within each paragraph. On a branch, there are seven birds: a raven, a hummingbird, a robin, a crow, a quail, a blue jay, and a cardinal. The hummingbird is to the left of the crow. The quail is to the right of the crow. The raven is to the left of the robin. The blue jay is the third from the right. The cardinal is the leftmost. The hummingbird is the fourth from the left.} \\
\texttt{Options:} \\
\texttt{(A) The raven is the rightmost} \\
\texttt{(B) The hummingbird is the rightmost} \\
\texttt{(C) The robin is the rightmost} \\
\texttt{(D) The crow is the rightmost} \\
\texttt{(E) The quail is the rightmost} \\
\texttt{(F) The blue jay is the rightmost} \\
\texttt{(G) The cardinal is the rightmost} \\
\texttt{A:} \\
            }
        }
    }
     \\ \\

    \theutterance \stepcounter{utterance}  
    & \multicolumn{4}{p{0.6\linewidth}}{
        \cellcolor[rgb]{0.561,0.737,0.561}{
            \makecell[{{p{\linewidth}}}]{
                \texttt{\tiny{[P1$\rangle$GM]}}
                \texttt{Let's analyze the given statements to determine the order of the birds:} \\
\\ 
\texttt{1. The hummingbird is to the left of the crow.} \\
\texttt{2. The quail is to the right of the crow.} \\
\texttt{3. The raven is to the left of the robin.} \\
\texttt{4. The blue jay is the third from the right.} \\
\texttt{5. The cardinal is the leftmost.} \\
\texttt{6. The hummingbird is the fourth from the left.} \\
\\ 
\texttt{From statement 5, we know the cardinal is in the first position.} \\
\texttt{From statement 6, we know the hummingbird is in the fourth position.} \\
\texttt{From statement 1, the crow must be to the right of the hummingbird, so it can be in positions 5, 6, or 7.} \\
\texttt{From statement 2, the quail must be to the right of the crow, so if the crow is in position 5, the quail can be in positions 6 or 7. If the crow is in position 6, the quail must be in position 7.} \\
\texttt{From statement 4, the blue jay is in the fifth position (third from the right in a sequence of seven).} \\
\texttt{From statement 3, the raven is to the left of the robin.} \\
\\ 
\texttt{Now let's place the remaining birds:} \\
\texttt{{-} Cardinal is in position 1.} \\
\texttt{{-} Hummingbird is in position 4.} \\
\texttt{{-} Blue jay is in position 5.} \\
\\ 
\texttt{This leaves positions 2, 3, 6, and 7 for the raven, robin, crow, and quail.} \\
\texttt{Since the crow must be to the right of the hummingbird (position 4), and the blue jay is in position 5, the crow can only be in position 6 or 7.} \\
\texttt{If the crow is in position 6, the quail must be in position 7 (since it's to the right of the crow).} \\
\texttt{This leaves positions 2 and 3 for the raven and robin. Since the raven is to the left of the robin, the raven must be in position 2 and the robin in position 3.} \\
\\ 
\texttt{So the order is:} \\
\texttt{1. Cardinal} \\
\texttt{2. Raven} \\
\texttt{3. Robin} \\
\texttt{4. Hummingbird} \\
\texttt{5. Blue jay} \\
\texttt{6. Crow} \\
\texttt{7. Quail} \\
\\ 
\texttt{Therefore, the rightmost bird is the quail.} \\
\\ 
\texttt{A: (E) The quail is the rightmost} \\
            }
        }
    }
    & & \\ \\

    \theutterance \stepcounter{utterance}  
    & & & \multicolumn{2}{p{0.3\linewidth}}{
        \cellcolor[rgb]{0.9,0.9,0.9}{
            \makecell[{{p{\linewidth}}}]{
                \texttt{\tiny{[GM$|$GM]}}
                \texttt{Target: (E)} \\
            }
        }
    }
    & & \\ \\

    \theutterance \stepcounter{utterance}  
    & & & \multicolumn{2}{p{0.3\linewidth}}{
        \cellcolor[rgb]{0.9,0.9,0.9}{
            \makecell[{{p{\linewidth}}}]{
                \texttt{\tiny{[GM$|$GM]}}
                \texttt{game\_result = LOSE} \\
            }
        }
    }
    & & \\ \\

\end{supertabular}
}

\end{document}
